\section{Related Work}
\label{sec:related}

\paragraph{Higher-level supercompilation.}
S{\o}rensen, Gl{\"u}ck, and Jones~\cite{sorensen-gluck-jones-1996-positive}
introduce \emph{positive supercompilation}, which goes beyond standard SC by
permitting generalisations that introduce new function symbols not present in
the original program.  This is the first instance of HLS.  Mitchell and
Runciman~\cite{mitchell-runciman-2010-supercompilation} implement HLS for
Haskell (Supero), showing that accumulator patterns and fusion identities can
be found automatically.  Bolingbroke~\cite{bolingbroke-2013-phd} develops
call-by-need SC with a rich generalisation algorithm.

Our contribution is orthogonal: we provide a \emph{proof-theoretic} account
of what HLS is doing (cut introduction), which leads directly to the LLM
oracle architecture.  Prior HLS work has not connected the generalisation
step to cut introduction, nor has it given a mechanised correctness proof.

\paragraph{Worker/wrapper.}
The worker/wrapper transformation~\cite{gill-hutton-2009-worker-wrapper}
is a structured form of accumulator introduction.  It can be seen as a special
case of cut introduction: the wrapper is the main program, the worker is the
cut formula.  HLS with LLM oracle can discover worker/wrapper transformations
automatically; mechanising this connection is future work.

\paragraph{Cyclic proofs.}
Brotherston and Simpson~\cite{brotherston-simpson-2011-sequent} establish
the sequent-calculus foundation for cyclic proofs.
Afshari and Wehr~\cite{afshari-wehr-2024-abstract-cyclic} give an abstract
categorical account.  Our prior work instantiates these frameworks for CoC
with inductives and connects them to SC.  The present paper uses that
connection as a design guide.

\paragraph{Formally verified supercompilation.}
Krustev~\cite{krustev-2014-simple-sc-coq,krustev-2018-sc-flia} gives the first
mechanised correctness proof of a supercompiler in Coq, verifying that the
residual program is observationally equivalent to the source.  His framework
handles a first-order language (Whistle) with homeomorphic embedding.
We extend this direction: (i)~our SC operates on CoC with inductives
(dependent types), (ii)~the CIU theorem lifts Krustev's observational
equivalence to the dependently-typed setting, and (iii)~the LLM oracle
and $\omega$-rule add no new proof obligations because the kernel
validation is inherited from the CIU theorem.

\paragraph{Bisimulation proofs for transformations.}
Jones and Hamilton~\cite{jones-hamilton-2006-bisimulation} prove the correctness
of unfold/fold transformations using bisimulation on program configurations.
Hamilton's later work on distillation~\cite{hamilton-2007-distillation} extends
this to labelled transition systems.  Our CIU theorem is a bisimulation proof
in the sense of~\cite{jones-hamilton-2006-bisimulation}: it establishes that
the original and residual terms are indistinguishable by any observational
context.  The proof-theoretic framing makes the generalisation step (fold)
explicit as cut introduction, rather than an implicit bisimulation up-to
technique.

\paragraph{LLMs in formal methods.}
Recent work uses LLMs to suggest proof steps in Lean and Isabelle.
Our usage is different: the LLM proposes \emph{program transformations}
(generalisation candidates), not proof terms.  The kernel check
(\texttt{trace\_condition\_ok}) plays the role of a type checker for
transformation proposals, rather than an oracle for proof search.
